\section{Resource characterizations and verification note}

The uniform comparison in \cref{obj:thm:uniform-annotation-frontier} supplies the common basis for the resource conclusions in \cref{obj:thm:annotation-resource-phases}. For every allowed experiment sequence in \cref{obj:synth_8}, the rate definition in \cref{obj:def:rate-handle} gives
\[
R(n_k,m_k,d_k,\epsilon_k)
\asymp
\min\left\{
1,\frac{1}{S_k}
+\left(\frac{d_k}{N_k\epsilon_k\ell_k}\right)^2
\right\}.
\]
Here the sequence-specific scales are those introduced in \cref{obj:def:phases}. The comparison constants are independent of the public indices. Consequently, the resource characterizations apply to arbitrary allowed sequences, including sequences with oscillating budgets, expanding alphabets, and diminishing overlap.

\subsection{Consistency and oracle order}

The consistency characterization in \cref{obj:thm:annotation-resource-phases} assigns a separate requirement to each component of the rate:
\[
\mathbf v\in\mathcal C
\quad\Longleftrightarrow\quad
S_k\longrightarrow\infty,
\qquad
\frac{d_k}{N_k\epsilon_k\ell_k}\longrightarrow0.
\]
The first condition concerns the information supplied by rare-arm outcomes in the complete records. The second concerns marginal information relative to the alphabet size. Auxiliary records contribute through the second denominator, while the complete-record budget and overlap floor determine the first scale. These two requirements characterize vanishing minimax risk throughout the sequence domain.

Oracle order compares precision with the rare-label benchmark in \cref{obj:prop:benchmark-recovery}. Under \(S_k\to\infty\), that benchmark eventually equals \(S_k^{-1}\). The oracle characterization therefore places the many-cell component at or below this inverse-information order:
\[
\mathbf v\in\mathcal O
\quad\Longleftrightarrow\quad
S_k\longrightarrow\infty,
\qquad
d_k=O\left(
\frac{N_k\epsilon_k\ell_k}{\sqrt{S_k}}
\right).
\]
Equivalently, the total marginal-information budget satisfies
\[
\frac{d_k\sqrt{S_k}}{\epsilon_k\ell_k}=O(N_k).
\]
As stated in \cref{obj:thm:annotation-resource-phases}, this budget characterization is conditional on divergence of the rare-arm information scale. Its order constant may depend on the sequence, whereas the underlying risk comparison is uniform over public indices.

For the supervised specialization \(m_k=0\), the identity \(N_k\epsilon_k=S_k\) gives the consistency threshold \(d_k=o(S_k\ell_k)\) and the oracle threshold \(d_k=O(\sqrt{S_k}\ell_k)\), together with \(S_k\to\infty\). These specializations identify the marginal-information requirements supplied entirely by complete records.

\subsection{Strict relative improvement}

The improvement phase in \cref{obj:def:phases} compares the auxiliary-record experiment with the supervised experiment at the same complete-record budget, alphabet size, and overlap floor. Applying the uniform risk comparison to both experiments gives
\[
\frac{R(n_k,m_k,d_k,\epsilon_k)}
     {R(n_k,0,d_k,\epsilon_k)}
\asymp
\frac{
\min\left\{1,S_k^{-1}
+\left(d_k/(N_k\epsilon_k\ell_k)\right)^2\right\}
}{
\min\left\{1,S_k^{-1}
+\left(d_k/(S_k\ell_k)\right)^2\right\}
}.
\]
Both caps enter this comparison, accommodating supervised risks ranging from vanishing precision to saturation.

The exact characterization in \cref{obj:thm:annotation-resource-phases} combines three requirements:
\[
S_k\longrightarrow\infty,
\qquad
\frac{d_k}{\sqrt{S_k}\ell_k}\longrightarrow\infty,
\qquad
\frac{N_k/n_k}
{\max\{1,d_k/(S_k\ell_k)\}}
\longrightarrow\infty.
\]
The first supports increasing outcome precision. The second places the supervised experiment above rare-label oracle order by a diverging factor. The third measures the marginal expansion relative to the larger of a constant and the supervised many-cell scale. Together they characterize a vanishing relative risk, including for sequences that move between different supervised resource regions. Within the supervised specialization, the relative risk ratio equals one.

\subsection{Saturation and the precision plateau}

The finite-experiment conclusions in \cref{obj:thm:annotation-resource-phases} retain the cap in the rate. When
\[
N\epsilon\ell\le d,
\]
the many-cell component reaches the cap, and the minimax risk is bounded above and below by positive universal constants. This saturation condition identifies a region of constant-order risk.

At the larger marginal budget
\[
N\ge \frac{d\sqrt{S}}{\epsilon\ell},
\]
the same result establishes
\[
R(n,m,d,\epsilon)\asymp b(n,\epsilon)
=\min\{1,S^{-1}\}.
\]
The cap preserves the comparison for small rare-arm information scales. Along every oracle sequence, \(S_k\to\infty\), so the risk eventually has order \(S_k^{-1}\).

The plateau conclusion also preserves the quantification over auxiliary budgets. By \cref{obj:thm:annotation-resource-phases}, on every oracle sequence there is a positive lower comparison constant such that, for all sufficiently large indices and every replacement budget \(m'\in\mathbb N\),
\[
R(n_k,m',d_k,\epsilon_k)\ge \frac{c}{S_k}.
\]
The rare-label floor in \cref{obj:thm:rare-label-floor} supplies this bound uniformly over the auxiliary channel. Thus oracle attainment identifies a precision order governed by complete-record outcome information, with further marginal information operating within that order.

\subsection{Verification note}

The mathematical scope includes the canonical randomized minimax risks, the explicit estimator and its risk guarantee, the testing and approximation bounds, the benchmark comparisons, the resource characterizations, and causal realization and identification. The statistical results use the finite observable model in \cref{obj:def:model}, the overlap and product-sampling conditions in \cref{obj:ass:overlap,obj:ass:data-product}, and the rule domains in \cref{obj:def:risk,obj:def:known-risk}. Causal identification concerns compatible potential-outcome extensions satisfying \cref{obj:ass:consistency,obj:ass:exchangeability}, as specified in \cref{obj:lem:causal-realization}. Auxiliary results retain their displayed independence, calibration, and approximation hypotheses. These conditions are assumed mathematical inputs to the stated implications.
